\documentclass[12pt]{article}

\usepackage[T1]{fontenc}
\usepackage[spanish]{babel}
\usepackage[letterpaper, margin=2.54cm]{geometry}
\usepackage{amsmath}

\usepackage[
	backend=biber,
	style=apa,
	]{biblatex}

\addbibresource{referencias.bib}

\begin{document}
	
%Portada

\begin{titlepage}
	\centering
	{\scshape\LARGE Universidad Aztlán \par}
	\vspace{1cm}
	{\scshape\Large Sede Playa del Carmen \par}
	\vspace{3cm}
	{\scshape\Huge Capital, Activo y Pasivo \par}
	\vspace{1cm}
	{\scshape\Large Trabajo de Investigación \par}
	\vspace{3cm}
	{\Large Jose Rosendo Camal Teh \par}
	\vfill
	{\Large Ingeniería en Sistemas Computacionales \par}
	\vspace{0.5cm}
	{\Large Segundo Cuatrimestre \par}
	\vspace{0.5cm}
	{\Large Contabilidad Financiera \par}
	\vfill
	{\Large Profa. Guadalupe Parada Valencia \par}
	\vfill
	{\Large 16 de marzo de 2026 \par}
	\vfill
	{\Large Playa del Carmen, Quintana Roo \par}
\end{titlepage}

\tableofcontents
\newpage

% Introducción
\section{Introducción}

La base conceptual y económica de la voz \emph{capital} se incorpora en la contabilidad de una manera estricta y fragmentada. Además, establece un equilibrio de la naturaleza de los \emph{activos}, los \emph{pasivos} y el \emph{capital contable} que cimenta los fundamentos para que en un futuro haya una comprensión de esta como herramienta para la toma de decisiones.

% Desarrollo
\section{Capital, Activo y Pasivo}
	
	% Capital
	\subsection{Capital}
	
		% Capital: concepto
		\subsubsection{Concepto}

		El concepto económico de \textbf{capital} representa todos aquellos conjuntos de activos y bienes económicos que posee una entidad para producir mayor riqueza. Dicho significado nos permite entrever que el capital está ligado a la idea del dinero, bienes, derechos y propiedades que posee una entidad (persona u organización) con el objetivo general de utilizarlos para la generación de patrimonio. Por lo tanto, es pertinente interpretar que tanto el medio como el fin son recursos cuantificables en una métrica monetaria. 

		En otras palabras, el capital expresa la noción de que los recursos son utilizados para generar más de sí mismos, por lo que deriva en la acumulación del capital. Sin embargo, en la perspectiva contable dicho término no se maneja en el sentido general, sino que dispone de tecnicismos que analizan los distintos matices del vocablo.

		En la contabilidad, el término capital se emplea principalmente para referirse a la diferencia aritmética de los bienes y derechos que posee una entidad y de lo que aquella debe. Para ello, se opta por la voz de \textbf{capital contable}. No se acostumbra recurrir a la palabra \emph{capital} a secas a pesar de que en el ámbito económico es relativamente clara y útil, debido a que dicha base conceptual se fragmenta en diversas acepciones técnicas como \emph{capital contable}, \emph{capital social}, \emph{capital financiero}, etc. Lo que vuelve frágil y ambiguo el uso de ``capital'' y obliga a emplear denominaciones más específicas.

		 Otro punto del que se debe tomar nota es el término de capital contable está más orientado a las entidades con fines de lucro, mientras que \textbf{patrimonio contable} posee casi el mismo significado con la sutileza del enfoque hacia entidades no lucrativas.

		% Capital: tipos capital y patrimonio contables
		\subsubsection{Tipos de capital y patrimonio contables}

		Como se ha mencionado, se emplea \textbf{capital contable} y \textbf{patrimonio contable} para designar la diferencia aritmética de los activos y pasivos según la naturaleza de lucro de la entidad.

		El capital contable de las entidades con fines lucrativos se clasifica en:
		\begin{itemize}
			\item \textbf{Capital contribuido}: Está integrado por las aportaciones de los propietarios (por ejemplo: dueños y accionistas) recibidos por la entidad y de otros instrumentos que la misma entidad ha emitido y que se cataloga como capital.
				\begin{itemize}
					\item Capital social.
					\item Aportaciones para futuros aumentos de capital.
					\item Prima de venta en acciones.
				\end{itemize}
			\item \textbf{Capital ganado}: Se conforma por las utilidades y pérdidas que ha acumulado la entidad a raíces de las aportaciones, deudas y de sus operaciones.
				\begin{itemize}
					\item Utilidades retenidas, incluyendo las aplicadas a reservas de capital.
					\item Pérdidas acumuladas.
					\item Exceso o suficiencia en la actualización del capital contable.
				\end{itemize}
		\end{itemize}

	% Activo
	\subsection{Activo}

		% Activo: concepto
		\subsubsection{Concepto}

		El activo es un recurso controlado y propiedad de una entidad del que espera obtener beneficios económicos en el futuro. Es un derivado de sus actividades u operaciones del pasado y es cuantificable en términos monetarios. Una entidad no suele tener un único activo sino un conjunto de ellos. Los activos de una entidad varían en naturaleza desde efectivo, mercancía y bonos del gobierno hasta patentes, marcas, edificios, maquinaria y transporte, por mencionar ejemplos concretos.

		% Activo: cuentas del activo
		\subsubsection{Principales cuentas del activo}

			Se listan las principales cuentas del activo: \\

			\textbf{Efectivo y bancos.} El efectivo está constituido por la moneda de curso legal y por los depósitos bancarios disponibles para la operación de la entidad.
			
			\textbf{Inversiones temporales.} Las inversiones temporales (o inversiones disponibles a la vista) son instrumentos de inversión convertibles en efectivo en el corto plazo y tienen el objeto de generar un rendimiento hasta el momento en que la entidad requiera liquidez.

			\textbf{Cuentas y documentos por cobrar.} Son las cuentas que representan derechos exigibles originados por ventas, servicios prestados, otorgamiento de préstamos o cualquier otro concepto análogo. Este grupo engloba cuentas de relevancia significativa como \emph{Clientes}, \emph{Documentos por cobrar}, \emph{Funcionarios y empleados}, \emph{Deudores diversos}, \emph{Intereses por cobrar}, \emph{Dividendos por cobrar}, \emph{Rentas por cobrar}, entre otros. De las cuentas mencionadas las más conocidas son clientes y documentos por cobrar.

			\textbf{Inventarios.} Son los activos que encajan en alguno de los tres criterios siguientes: (1) tiene la finalidad de ser vendido; (2) estar en proceso de producción para luego ser vendido; y (3) ser utilizado como suministro, material o materia prima en el proceso de producción o en la prestación de un servicio.

			\textbf{Propiedades, plantas y equipo.} Este grupo reúne las cuentas de \emph{Propiedades}, \emph{Maquinaria}, \emph{Terrenos}, \emph{Edificios}, \emph{Equipo de oficina}, entre otras. Son activos tangibles de la entidad que tiene como finalidad el usarlos en la producción o en el suministro de bienes y servicios, en la prestación de servicios a la entidad, a sus clientes o al público en general. En un plazo, por lo general, mayor a un año.

			\textbf{Activos intangibles}. Este concepto representa a los activos que carecen de sustancia física, identificables y con capacidad de generar beneficios a la entidad. El acceso a estos beneficios pueden ser controlados por la entidad a través de mecanismos legales para evitar que terceros disfruten de estos frutos. Las cuentas más importantes son \emph{Patentes}, \emph{Derechos de autor}, \emph{Fórmulas secretas}, \emph{Licencias}, \emph{Gastos en Investigación y Desarrollo}, entre otras. 

		% Activo: clasificación
		\subsubsection{Clasificación}

			Los activos se clasifican por su grado de disponibilidad y esto es la facilidad de conversión o transformación en efectivo. Los recursos con mayor grado de disponibilidad son los \emph{activos circulantes}, mientras que los de menor disponibilidad son aquellos que no son fácilmente convertibles en efectivo (\emph{activos no circulantes}). \\

			\textbf{Activo circulante.} Este grupo de activos está conformado por los activos cuyos beneficios económicos son obtenidos en un ciclo de operación normal o un periodo no mayor a un año. Son fácilmente convertibles a efectivo, suelen estar en movimiento y rotación constantes, y fueron adquiridos con el propósito de venderlos.

			Como ilustración, algunas de las principales cuentas del activo circulante son:

			\begin{itemize}
				\item \textbf{Efectivo}
				\item \textbf{Bancos}
				\item \textbf{Inversiones temporales}
				\item \textbf{Clientes}
				\item \textbf{Cuentas y documentos por cobrar}
				\item \textbf{Inventarios}
			\end{itemize}

			\textbf{Activo no circulante.} Los activos que conforman a este grupo son aquellos del que sus beneficios son obtenidos a un plazo mayor a un ciclo de operación normal o a un periodo mayor de un año. No son fácilmente convertibles a efectivo y por lo general, son activos fijos y permanentes. El propósito de su adquisición es emplearlos en sus operaciones y no en venderlos.

			Ejemplos de principales activos que conforman este conjunto son:

			\begin{itemize}
				\item \textbf{Equipo de transporte}
				\item \textbf{Maquinaria}
				\item \textbf{Terrenos}
				\item \textbf{Edificios}
				\item \textbf{Equipo de oficina}
				\item \textbf{Patentes, marcas registradas}
			\end{itemize}

	% Pasivo
	\subsection{Pasivo}

		% Pasivo: concepto
		\subsubsection{Concepto}

		El pasivo es una deuda o una obligación contraída por una entidad del que espera obtener pérdidas económicas o una disminución de los beneficios económicos en el futuro. Son consecuencias de sus actividades u operaciones del pasado y es cuantificable en términos monetarios. Las entidades no poseen un solo pasivo sino un conjunto de ellos y esto no es necesariamente una señal de alarma. Un pasivo no es por sí solo un indicador negativo puesto que su naturaleza varía desde préstamos y créditos hasta pagos de sueldos y salarios, comisiones y diversos impuestos. Los pasivos sin un debido control pueden convertirse en una carga económica y eso depende más de la gestión que les otorgue la entidad que de los propios pasivos.
			
		% Pasivo: cuentas del pasivo
		\subsubsection{Principales cuentas del pasivo}

			Se presentan las principales cuentas del pasivo: \\
			
			\textbf{Proveedores.} Esta cuenta representa los importes y adeudos que una entidad debe a otras entidades por la adquisición de mercancías u otros conceptos que se pagan en un periodo normal de operaciones, dentro de un año o menos.

			\textbf{Documentos por pagar.} Son los títulos de crédito a cargo de la entidad como pagarés o documentos con interés descontado (recibe menos de lo que se otorgó en crédito).

			\textbf{Acreedores diversos.} Son las entidades a las que una entidad les debe por conceptos distintos a la adquisición de mercancías.

			\textbf{Impuestos por pagar.} Son los impuestos locales y federales que una entidad debe pagar que derivan de la naturaleza de su operación. Los impuestos de carácter fiscal como el \emph{ISR} y el \emph{IVA}, así como los de origen laboral y de seguridad social, como las cuotas del \emph{IMSS}, del \emph{SAR} y del \emph{INFONAVIT}, por mencionar algunos ejemplos.

			\textbf{Gastos por pagar.} Son los gastos por servicios o beneficios obtenidos que aún están pendientes de pagarse. Algunos ejemplos son los \emph{Sueldos y salarios}, \emph{Comisiones}, \emph{Servicios básicos}, \emph{Arrendamiento}, etc.

			\textbf{Anticipo de clientes.} Son los activos de clientes que aseguran las futuras ventas de mercancías o prestación de servicios de una entidad con la finalidad de esta cuente con los recursos necesarios para la producción o prestación de bienes y servicios.

			\textbf{Hipotecas por pagar.} Esta cuenta agrupa a las transacciones de las que se obtienen recursos de una institución bancaria con la promesa de pagarlos a largo plazo cuya garantía es la escritura de edificios, departamentos o terrenos.


		% Pasivo: clasificación
		\subsubsection{Clasificación}

			Los pasivos se clasifican por su grado de exigibilidad. Esto se debe entender como el periodo de vencimiento y la urgencia de liquidación de las deudas y obligaciones adquiridos por la entidad. Las obligaciones y los compromisos de mayor grado de exigibilidad son los \emph{pasivos a corto plazo} y lo de menor grado de exigibilidad son los \emph{pasivos a largo plazo}. \\

			\textbf{Pasivo a corto plazo.} Integrado por las deudas y obligaciones a cargo de la entidad cuyos vencimientos corresponden a un periodo menor de un año o a un ciclo de operación normal. Estos pasivos tienen la particularidad de tener un mayor grado de exigibilidad, es decir, su periodo de vencimiento es breve [como se ha mencionado] y hay una mayor presión para su liquidación.
				
			Algunas de las principales deudas y obligaciones que constituyen este grupo son:

			\begin{itemize}
				\item \textbf{Proveedores}
				\item \textbf{Documentos por pagar}
				\item \textbf{Acreedores diversos}
				\item \textbf{Impuestos por pagar}
				\item \textbf{Gastos por pagar}
				\item \textbf{Anticipo de clientes}
				\item \textbf{Dividendos por pagar}
			\end{itemize}

			\textbf{Pasivo a largo plazo.} Este grupo está conformado por las obligaciones y los compromisos a cargo de la entidad del que no tiene mucha preocupación de liquidarlos en un breve periodo de tiempo. En otras palabras, tiene un periodo de vencimiento mayor a un ciclo de operación normal o mayor a un año.

			A continuación se listan algunas de las principales deudas y obligaciones que constituyen este grupo:

			\begin{itemize}
				\item \textbf{Documentos por pagar a largo plazo}
				\item \textbf{Obligaciones en circulación}
				\item \textbf{Hipotecas por pagar}
				\item \textbf{Préstamos bancarios}
			\end{itemize}

	% Relación entre capital, activo y pasivo
	\subsection{Relación entre capital, activo y pasivo}

	El núcleo de la contabilidad es la \textbf{ecuación contable}, la cual representa la relación entre los conceptos de capital contable, activo y pasivo. En esta disciplina, los fundamentos recaen en la base conceptual de las concepciones de \emph{capital (económico)} y del equilibrio de los componentes de la \emph{ecuación contable}.

	La ecuación contable indica en el miembro izquierdo los recursos económicos que posee y dispone la entidad (\textbf{activos}) y en el miembro derecho las obligaciones, compromisos y fuentes de financiamiento (\textbf{pasivos} y \textbf{capital contable}). Esta relación se formaliza como sigue:

	\begin{align*}
		\textit{Activos} &= \textit{Pasivos} + \textit{Capital contable}
	\end{align*}

	El activo, el pasivo y el capital contable se desglosan de esta forma:

	\begin{align*}
		\textit{Activos} &= \textit{Activos circulantes} + \textit{Activos no circulantes} \\
		\textit{Pasivos} &= \textit{Pasivos a corto plazo} + \textit{Pasivos a largo plazo} \\
		\textit{Capital Contable} &= \textit{Capital contribuido} + \textit{Capital ganado}
	\end{align*}	

	Bajo la perspectiva matemática podemos despejar el activo, el pasivo y el capital contable y las cuentas que los constituyen. En otros términos, esta relación se condensa en una operación de índole aritmética y algebraica en donde la contabilidad les brinda un marco teórico cuya aplicación es útil para las entidades (personas, empresas, instituciones y gobiernos) para la toma de decisiones.

% Conclusión
\section{Conclusión}

El \emph{activo}, el \emph{pasivo} y el \emph{capital contable} son tecnicismos de la rama contable cuyas esencias emergen de la concepción de capital en el ámbito económico. De manera general, estas tres voces se enfocan en sintetizar y clasificar los bienes, derechos, obligaciones y compromisos bajo un marco contable cuya formalización es la \emph{ecuación contable}. Esto asienta el antecedente para más tarde visualizar la contabilidad como un instrumento de toma de decisiones.

% Referencias

\nocite{*}
\printbibliography[heading=bibintoc]

\end{document}
